We prove Serre duality. For now only for vector bundles.



\subsection{Duality for Projective Space}

Given a bundle $\mc F : \bP^n \to \Mod R$ there is a natural pairing
\[
  \Hom_X(\mc F,\omega) \times H^n(X,\mc F) \to H^n(X,\omega), \;
  (\varphi,\xi) \mapsto H^n(\varphi)(\xi)
  \rlap{.}
\]
Since $H^n(\bP^n,\omega) = R$ by \Cref{canonical-bundle-as-twist}, it induces a map
\begin{align*}
  \vartheta^0(\mc F) : \Hom_X(\mc F,\omega) \to H^n(X,\mc F)^\ast \\
\end{align*}
where $X \equiv \bP^n$.

\begin{proposition}
  If $\mc F$ is finitely presented, then $\vartheta^0(\mc F)$ is an isomorphism.
\end{proposition}

\begin{proof}
  In the case $\mathcal F\simeq \mathcal O(q)$ we compute $\Hom_X(\mathcal O(q),\omega)\simeq H^0(\bP^n,\omega\otimes \mathcal O(-q))$.
  So in this case, we can conclude as we know the cohomology of twisting sheaves.
  By \Cref{fp-module-bundle-globally-generated}, we get an exact sequence:
  \[ \mathcal O(p)^{\oplus n}\to\mathcal O(p')^{\oplus n'}\to \mathcal F\to 0 \]
  and the functors $\Hom_X(-,\omega)$ and $H^n(X,-)^\ast$ both map right exact sequences to left exact sequences. So we can conclude with the 5-lemma.
\end{proof}

\begin{proposition}
  If $\mathcal E : X \to \Mod R$ is a vector bundle then for every $i\geq 0$ there are natural isomorphisms
  \[
    \vartheta^i(\mathcal E) : H^i(X, \mathcal E^\ast \otimes \omega) \to H^{n-i}(X,\mathcal E)^\ast
  \]
  such that $\vartheta^0(\mc E)$ agrees with the corresponding map above.
\end{proposition}

\begin{proof}
  Since vector bundles have finitely adequate cohomology by \Cref{vector-bundles-have-faq-cohomology-which-eventually-vanishes} we know that both $H^i(X, (-)^\ast \otimes \omega)$ and $H^{n-i}(X,-)^\ast$ are contravariant and exact $\delta$-functors.
    These are also both coherently coeffaceable by \Cref{vector-bundles-can-be-coherently-resolved} and thus both universal by \Cref{coherently-effaceable-implies-universal}.
\end{proof}

\begin{remark}
  Already on $\bP^1$, the map $\vartheta^0(\mc F)$ need not be an isomorphism for finitely copresented bundles.
  In other words, the above pairing is not perfect for finitely copresented bundles (or -- by a dual argument to the one we are about to state -- for finitely presented bundles).
  Indeed, let $\mc K, \mc C$ and $\mc I$ be respectively the kernel, cokernel and image of the map
  \[
    \mc O(-2) \overset{\cdot x}{\longrightarrow} \mc O(-1)
  \]
  on $\bP^1$.
  Thus $\mc C(1) = \mc K^\ast \otimes \mc O(-2)$.
  We will show that $H^0(\bP^1,\mc C(1))$ and $H^1(\bP^1,\mc K)^\ast$ are not isomorphic, violating Serre-duality.
  Note that while $0 \to \mc O(-2) \to \mc O(-1) \to \mc C \to 0$ is not pointwise exact, it is affine-locally exact.
  Thus, using Čech cohomology, we see that $\mc K$ has no cohomology at all, in particular $H^1(\bP^1,\mc K)^\ast = 0$.
  Similarly we see $H^0(\bP^1, \mc C(1))=R$, thus finishing the argument.
\end{remark}

\begin{remark}
  On $\bP^2$, Serre duality -- at least when formulated as above -- fails even for finitely presented bundles.
  To see this, we define a finitely presented bundle $\mc C$ as the cokernel
  \[
    \mc O(-3) \overset{(x_0,x_1)}{\longrightarrow} \mc O(-2)^{\oplus 2} \longrightarrow \mc C \to 0 \rlap{.}
  \]
  Using a similar Čech cohomology calculation as in the previous remark {\color{red}TODO: Details}, we see that $H^1(\bP^2,\mc C^\ast \otimes \omega_{\bP^2}) = 0$ and $H^1(\bP^2, \mc C)^\ast = R$.
  In particular, $\vartheta^1(\mc C)$ cannot be an isomorphism.
\end{remark}

Note that the above example works even in classical algebraic geometry.
If one wants to generalize the statement to a bigger class of bundles, then one needs to replace $H^i(X,\mc F^\ast \otimes \omega_X)$ by a suitable notion of $\Ext^i(\mc F,\omega_X)$.
Roughly, the problem is as follows.

If one sets $\Ext^i(\mc F,\mc G) :\equiv H^i(X, \uHom(\mc F, \mc G))$ -- as was silently done in the case of vector bundles -- then the resulting spaces only classify pointwise split extensions.
This is exactly the entrypoint of the counterexamples of the above remarks:
The appearing sequences are not pointwise exact, hence do not constitute an element of $\Ext^1$, but are affine-locally exact.
In both examples, the $\Ext^1$-modules vanish, while the corresponding modules, that Serre duality wants them to compare to, do not.
In this sense, they are too small because they classify a too restictive notion of exactness.



